\section{A prescribed source and a terminal interval}\label{c:window}
Prescribing one endpoint and summing the other over a separated interval
gives normalizing mass of order $R^{-1/4}$. We first establish this
normalization by constructing interior connectors. We then prove the
bridge and attachment estimates that transfer the length law to these
chords. The attachment estimates will also be used for nonparallel cuts
in Section~\ref{c:tilted}.

\subsection{Interior connectors}

\begin{proposition}[Interior connectors]\label{c:connectors}
For two macroscopically separated side-interior ports in the hexagons of
Theorem~\ref{c:main}, the critical mass of connecting paths inside the
hexagon is at least $cR^{-5/4}$. The same bound holds in a parallel strip
of thickness comparable to $R$, with one port on each opposite cut and
displacement $O(R)$, and in compact nondegenerate wedge geometries with
one port on each ray, as in Theorem~\ref{c:main}. The constants are uniform
on fixed compact sets of these scaled configurations.
\end{proposition}
\begin{proof}
We adapt the cut-averaging construction in the proof of
\cite[Proposition 10.2]{UniformMarked}, placing the routes in the interior
rather than the exterior. The geometric choices below ensure that the
decomposition can be recovered for two different lists of cuts.

\paragraph{Routes and junction cuts.}
Work first in a hexagon and scale distances by $R$. Start with short
segments in the inward normals at the two ports. Continue by disjoint
simple arcs in the open interior to two parallel lanes in a free interior
ball, preserving their cyclic order. The lanes approach one common pure
transverse line from the same side. Their separation is a small fixed
positive number, so a completing arch of diameter at most a fixed
multiple of that separation fits in the free half-ball beyond the line.
Keep that half-ball disjoint from all nonincident route portions.

Approximate the arcs by finitely many segments in pure normal directions,
retaining the initial and terminal straight portions and making successive
direction changes at most $\pi/3$. One can do this by alternating the two
adjacent pure directions spanning each smooth tangent, with a short step
in their common direction at a sector transition. A sufficiently fine
fixed mesh preserves separation of nonlocal portions; on nearby portions,
positive projection in a common direction preserves simplicity. Every
segment has length a fixed positive multiple of $R$.

Index a branch's segments by their guiding directions $n_j$. The cut at
the end of segment $j$ has normal $n_j$, the \emph{arriving} direction.
The next leg starts with inward normal $n_j$ and travels toward a cut
with outward normal $n_{j+1}$. Proposition~\ref{c:finite}(iii) therefore
keeps the incoming leg before the junction cut and the outgoing leg
after it. Restrict each leg to a sufficiently narrow corridor about its
segment. Each junction has its own small crossing neighborhood, avoided
by all nonincident route portions. This is a local requirement: those
portions need not avoid the entire infinite line of the cut.

The geometric choices have positive clearances. They persist in a
neighborhood of the scaled endpoint data, including the small support
perturbations allowed in Theorem~\ref{c:main}. Compactness gives uniform
choices of the number of segments and their relative tolerances.

\paragraph{One list and two lists of levels.}
Vary each junction cut in a fixed small window of admissible levels of
width comparable to $R$. Choose the windows so that all corridors and
stage orders just described persist for every list. The terminal level
is shared by the two branches; all other levels vary independently.
If there are $k$ legs in total, the set $\Lambda_R$ of level lists is
thus a product of $k-2$ ordinary windows and one shared terminal window.
Give it the uniform product probability measure.

For a list $\lambda$, let $I_\lambda(\gamma)$ indicate that the chord
splits into the confined legs and the completing arch at these cuts.
The splitting is unique. At each isolated junction the two incident
pieces meet the cut once and occupy its opposite half-planes; all other
pieces avoid that crossing neighborhood. Put
\[
 F_R(\gamma)=\frac1{|\Lambda_R|}\sum_{\lambda\in\Lambda_R}I_\lambda(\gamma),
 \qquad
 \int G\,d\mu=\sum_{\gamma:a\to b}\kappa^{L(\gamma)}G(\gamma).
\]
Multiplication of the $k$ point-to-line lower bounds and the confined
arch lower bound gives
\[
 \int F_R\,d\mu\ge cR^{-k/4-5/4}.
\]

For two lists, split each branch at both crossings in each junction
neighborhood. The two crossings occur in their normal-level order:
the incoming portions start below both levels, and the outgoing
portions end above both. The intervening portion is consequently a
strict bridge of the level difference. The portions between successive
neighborhoods still have separation comparable to $R$ and cost
$O(R^{-1/4})$ each. There are $k$ such long portions. Beyond the later
shared terminal line the remaining arch costs $O(R^{-5/4})$.

For two independent uniform levels $u,u'$ in a window of order $R$,
\[
 \mathbb E(1+|u-u'|)^{-\alpha}\le C_\alpha R^{-\alpha}
 \qquad(0<\alpha<1).
\]
Every ordinary junction gives one difference bridge and uses this bound
with $\alpha=1/4$. The shared terminal level gives two difference
bridges, one on each branch, and uses $\alpha=1/2$. Ties are included
by the convention $B_0=1$. Averaging the joint mass therefore gives
\[
 \int F_R^2\,d\mu
 \le C\underbrace{R^{-k/4}}_{\text{long portions}}
       \underbrace{R^{-(k-2)/4}R^{-1/2}}_{\text{difference bridges}}
       \underbrace{R^{-5/4}}_{\text{terminal arch}}
 =CR^{-2k/4-5/4}.
\]
Cauchy--Schwarz now yields
\[
 \mu(F_R>0)\ge
 \frac{(\int F_R\,d\mu)^2}{\int F_R^2\,d\mu}
 \ge cR^{-5/4}.
\]

For a wedge, add distant pure cuts at distance $AR$ from the endpoints,
with $A$ fixed and sufficiently large. The endpoint clearances to the
opposite supports permit initial inward segments of length $cR$.
The resulting bounded convex patch contains an interior ball of radius
$c'R$ accessible by the two separated routes above. After scaling, these
choices persist on a neighborhood of each nondegenerate configuration;
a finite cover of the compact family gives uniform clearances. The
constructed paths stay in the full wedge. A parallel strip admits the
same bounded-patch construction. This proves all the stated cases.
\end{proof}

\begin{theorem}[A source and a macroscopic terminal interval]\label{c:window-law}
In the hexagons of Theorem~\ref{c:main}, fix a source port tending after
scaling to a side-interior point. Let the terminal ports form a
macroscopic interval converging to a compact side-interior interval at
positive distance from the source limit. Under critical weights,
normalized over those chords, length is $R^{4/3+o(1)}$ in probability.
The same holds when the source is summed over a separated macroscopic
side-interior interval. For each fixed source the normalizer is
comparable to $R^{-1/4}$.
\end{theorem}

The diameter bound gives the upper normalization, and summing
Proposition~\ref{c:connectors} over order-$R$ terminal ports gives the
lower one. We prepare the length estimates next; the proof of the theorem
is completed at the end of this section.

\subsection{Bridge length and displacement windows}

\begin{lemma}[Fixed-height bridge length]\label{c:fixed-height-window}
Under critical weights on strict bridges of integral height $h$, with
one initial port fixed and the terminal port summed, for every $\delta>0$,
\[
 \Pr\{h^{4/3-\delta}\le L\le h^{4/3+\delta}\}\longrightarrow1.
\]
The same conclusion holds after any restriction retaining a fixed positive
fraction of $B_h$.
\end{lemma}
The fixed-height conclusion also appears in
\cite[Proposition 5.7]{Renewal}. We give the argument from the finite
and renewal inputs selected in Section~\ref{c:inputs}.
\begin{proof}
It suffices to treat $0<\delta<1$. The upper length tail follows from
$B_h^{-1}\sum Lw\le Ch^{4/3}$. For the lower tail, first obtain a
height window under the one-piece law $p$. There are fixed $M,A>1$ with
\begin{equation}\label{c:one-piece-height-window}
 p\{x\le H\le Mx,\ D\le Ax\}\ge cx^{-3/4}
\end{equation}
for all sufficiently large $x$. Indeed
$\mathbb E_p[1-e^{-H/(Kx)}]\asymp(Kx)^{-3/4}$, whereas the contribution
from $H\le x$ is $O(x^{-3/4}/K)$ and that from $H>Mx$ is
$O((Mx)^{-3/4})$. Choose first $K$, then $M$, sufficiently large.
Finally the diameter tail permits $A\ge M$ large enough to retain the
lower bound.

This window gives the concentration estimate
\begin{equation}\label{c:height-sum-atom}
 \sup_j\Pr(H_1+\cdots+H_k=j)\le Ck^{-4/3}.
\end{equation}
To see this, symmetrize the characteristic function using an independent
copy of $H$ that is bounded with positive probability. In
\eqref{c:one-piece-height-window}, take $x$ to be a sufficiently small
fixed multiple of $|\theta|^{-1}$. The cosine deficit on that event is
bounded below, so the squared modulus loses at least
$c|\theta|^{3/4}$ near zero. The modulus has only finitely many peaks
on the lattice torus, repeats the same bound near each, and is bounded
away from one on the remaining compact set. Fourier inversion proves
\eqref{c:height-sum-atom}.

For $k\ge2$, split the pieces into comparable independent halves.
If their total height is $h$, one half has height at least $h/2$, at
cost at most $Ckh^{-3/4}$ by the tail and truncated first moment.
The other half matches a height atom at cost $Ck^{-4/3}$. The case
$k=1$ costs at most $Ch^{-3/4}$. Hence, for fixed $\eta>0$,
\[
 \sum_{1\le k\le h^{3/4-\eta}}
 \Pr\Bigl\{\sum_{j\le k}H_j=h\Bigr\}
 \le Ch^{-3/4}\left(1+
       \sum_{2\le k\le h^{3/4-\eta}}k^{-1/3}\right)
 =o(B_h).
\]
For the remaining counts put $t=h^{4/3-\delta}$. Exponential Markov
and the length deficit give
\[
 \Pr\Bigl\{\sum_{j\le k}L_j<t\Bigr\}
 \le e(\mathbb E_p e^{-L/t})^k
 \le e\exp\{-k t^{-9/16+o(1)}\}.
\]
Choose $\eta<9\delta/16$. Summing this geometric bound over
$k>h^{3/4-\eta}$ is stretched-exponentially small. The renewal
factorization proves the lower length tail after division by $B_h$.
Dividing instead by any fixed positive fraction of $B_h$ preserves
both negligible exceptional masses.
\end{proof}

We next arrange a displacement that points inward at both chord ends.
This is needed when a translated chord is attached back to its original
supports. Distances in the following geometric statement are physical;
the bridge height $h$ is still in row units.

\begin{lemma}[Positive projections at two supports]\label{c:positive-projection}
Let $n,m$ be pure unit normals with $m\ne-n$. There are constants
$b,B,A,c_0>0$ such that, for every fixed $\delta>0$ and all sufficiently
large integral $h$, the mass of strict $n$-bridges from a fixed port
with endpoint displacement $v$, restricted by
\[
 bh\le m\cdot v\le Bh,\qquad D\le Ah,\qquad
 h^{4/3-\delta}\le L\le h^{4/3+\delta},
\]
is at least $c_0h^{-1/4}$. The constants may depend on $n,m$ but not
on $\delta$; the threshold may depend on $\delta$.
\end{lemma}
\begin{proof}
If $m\cdot n>0$, use a sufficiently narrow straight channel along $n$.
The remaining case is an angle of $2\pi/3$. Let $n'$ be the pure normal
at angle $\pi/3$ to both, and put $T=dh$. First travel to an $n'$-normal
cut at distance near $3T/2$ from the source, then travel in direction
$n$ to the top cut of the $n$-strip. Proposition~\ref{c:finite}(iii)
guides the first leg along $n'$, even though its initial inward normal
is $n$. At the intermediate position, both the $n$- and $m$-projections
are near $3T/4$. The final leg advances about $T/4$ along $n$, leaving
$m$-projection near $5T/8$. Small fixed tolerances keep both legs
between the $n$-strip boundaries, on opposite sides of the intermediate
cut, and within diameter $O(h)$.

Average that cut over a small window of order $h$. For one cut the two
legs have product mass at least $ch^{-1/2}$. For two cuts, recovery in
level order gives two long turned crossings and a difference bridge;
the averaged joint mass is at most $Ch^{-3/4}$. Cauchy--Schwarz gives
mass at least $ch^{-1/4}$ with the stated geometric restrictions.
In both angular cases this is a fixed fraction of $B_h$.
Lemma~\ref{c:fixed-height-window}, used with slightly smaller exponent
slack, removes the atypical lengths without changing the lower constant
for all sufficiently large $h$.
\end{proof}

\subsection{The cost of two translation choices}

Fix pure normals $n,m$ with $m\ne-n$. For a lattice translation $z$
with $n\cdot z>0$, let $D_n(z)$ be the critical mass of strict
$n$-bridges between the prescribed positions separated by $z$.
For $n\cdot z<0$, set $D_n(z)=D_n(-z)$; thus the bridge is always
oriented in increasing $n$-level order. At zero height set
$D_n(0)=1$ and $D_n(z)=0$ for nonzero $z$ with $n\cdot z=0$.

\begin{lemma}[Mixed displacement estimate]\label{c:mixed-displacement}
For every fixed $C_0<\infty$ there is $C<\infty$ such that, for $s\ge1$,
\begin{equation}\label{c:mixed-displacement-bound}
 \sum_{|z|\le C_0s}D_n(z)(1+|m\cdot z|)^{-1/4}\le Cs^{1/2}.
\end{equation}
If $Y$ is the $m$-projection of one $n$-directed irreducible displacement,
then its independent sums also satisfy
\begin{equation}\label{c:projected-concentration}
 \sup_y\Pr\Bigl\{\sum_{j\le k}Y_j=y\Bigr\}\le Ck^{-4/3},
 \qquad
 \sup_w\mathbb E\Bigl(1+\Bigl|w+\sum_{j\le k}Y_j\Bigr|\Bigr)^{-1/4}
 \le Ck^{-1/3}.
\end{equation}
\end{lemma}
\begin{proof}
The height-and-diameter window \eqref{c:one-piece-height-window} and
lateral reflection symmetry imply, for fixed $B>b>0$,
\[
 p\{bx\le |Y|\le Bx\}\ge cx^{-3/4}.
\]
Indeed the two reflected displacements have $m$-projections whose sum
is $2(m\cdot n)dH$. Since $m\cdot n\ne0$, at least one has magnitude
comparable to $H$, while the diameter cutoff bounds both from above.
The same symmetrization and Fourier argument as for
\eqref{c:height-sum-atom} gives the first bound in
\eqref{c:projected-concentration} on the projection's actual lattice.
For the second, sum this atom bound with weight $(1+|y|)^{-1/4}$
over $|y|\le k^{4/3}$, and bound the weight outside by $k^{-1/3}$.
This is uniform over translated lattices and the shift $w$.

It suffices to sum over $n\cdot z\ge0$ in
\eqref{c:mixed-displacement-bound}. Use renewal factorization, splitting
$k\ge2$ pieces into comparable halves. The displacement restriction
implies total height at most $C_1s$. Height damping bounds this event
on the first half by
\[
 C\exp\{-ck/s^{3/4}\}.
\]
Conditioning on that half, the second bound in
\eqref{c:projected-concentration} applies to the other half, with a
shift given by the first half's projection. Thus the sum is at most
\[
 C+C\sum_{k\ge2}e^{-ck/s^{3/4}}k^{-1/3}\le Cs^{1/2}.
\]
The constant covers one piece and the zero-height match. No coordinates
within one piece have been made independent.
\end{proof}

\subsection{Recoverable attachments}

Consider the two inward support normals $n,m$ at a source $a$ and a
terminal interval whose distances from $a$ are comparable to $R$.
When the supports differ, assume the endpoints have order-$R$ clearance
to the other support, as in the geometries above. We treat $m\ne-n$;
opposite supports will use Lemma~\ref{c:fixed-height-window} directly.
For identical supports their common half-plane replaces the two-cut
domain.

Choose a small fixed $c>0$ and a large scale $s\le cR$. Translate an
input chord by a lattice vector $v$ with $|v|\le Cs$ and both
\[
 g_1=n\cdot v/d,\qquad g_2=m\cdot v/d
\]
between fixed positive multiples of $s$. Attach a height-$g_i$ bridge
from each translated endpoint out to its original support. Restrict the
two bridges to diameter at most $C's$ and length between
$s^{4/3-\delta}$ and $s^{4/3+\delta}$. Such restrictions retain fixed
fractions of the available bridge mass, by the preceding lemmas and the
diameter tail. The branches stay in separated endpoint neighborhoods.
They avoid the translated chord by its supports, and, for distinct
supports, each stays beyond the translated support at the other end
by endpoint clearance and the small choice of $c$.

\begin{proposition}[Attachment moments]\label{c:attachment-proposition}
Let $\mathcal E$ be an event of input chords from $a$ to the stated
terminal interval, of critical mass $E$. The fixed-source attachment
construction just described has a translation count $N_s$ on output
chords and critical output measure $\mu_a$ such that
\begin{equation}\label{c:attachment-moments}
 \int N_s\,d\mu_a\ge cs^{1/2}E,
 \qquad
 \int N_s^2\,d\mu_a\le CsR^{-1/4}.
\end{equation}
Here the side-1 outer landing is exactly $a$; the side-2 outer landing
is free on its support. The estimate applies also to input chords in
the full two-cut domain, without a diameter bound.

If the input chords lie in a convex pure domain of diameter $O(R)$,
there is also a construction at scale $s=cR$ with both outer landings
free, all outputs in one convex pure domain of diameter $O(R)$, and
\begin{equation}\label{c:free-attachment-moments}
 \int N_R^{\rm free}\,d\mu_{\rm free}\ge cR^{3/2}E,
 \qquad
 \int(N_R^{\rm free})^2\,d\mu_{\rm free}\le CR^{7/4}.
\end{equation}
The measure $\mu_{\rm free}$ sums critical weights over ordered output
endpoint pairs. In both constructions output length is the input length
plus the two attached lengths. Constants are uniform for fixed compact
endpoint geometries and independent of the fixed length slack $\delta$;
the required lower threshold for $s$ may depend on $\delta$.
\end{proposition}
\begin{proof}
For fixed source $a$, reverse the side-1 bridge: it starts at $a$ and
must have displacement exactly $v$. Let $\mathcal V_s$ be all translations
in the size and positive-projection ranges above, with ranges wide enough
to include those in Lemma~\ref{c:positive-projection}. At each of order
$s$ heights, that lemma supplies restricted bridge mass at least
$c_0s^{-1/4}$ after summing over its allowed displacements. Thus the
total side-1 mass summed over $v$ is at least $cs^{3/4}$. Side 2 has
its outer landing free and supplies at least $cs^{-1/4}$ for each $v$.

Let $\mathcal A_v(\mathcal E)$ be the set of output chords obtained with
translation $v$, and define
\[
 N_s(\Gamma)=\sum_{v\in\mathcal V_s}
              \mathbf1_{\{\Gamma\in\mathcal A_v(\mathcal E)\}}.
\]
For a fixed $v$ the output determines its decomposition: the crossings
of the two translated supports separate the outer bridges from the
central translated chord. On identical supports there are two crossings,
localized near the two respective endpoints. Critical weights multiply,
since the joining ports have no weight. This proves the first inequality
in \eqref{c:attachment-moments}.

\paragraph{Recovering two translations.}
For $v,v'$, order the two crossings on each branch by that branch's
inward normal height. Their orders may disagree between the two ends.
Nevertheless both side-1 crossings occur before either side-2 crossing
along the output chord: the outer prefixes and suffixes are localized
in the two disjoint endpoint neighborhoods. Split into five portions:
the side-1 outer bridge, its difference bridge, a central chord, the
side-2 difference bridge, and the side-2 outer bridge.

The central portion is between the two innermost crossings. It stays
inside both corresponding supports, because each construction remains
beyond its own translated support after crossing, except for its
localized branch at the other end when the supports coincide. Those
other-end branches lie outside the central portion. Its initial point
is prescribed once the two translations are fixed, its terminal point
is free, and their separation is of order $R$. Dropping all other
conditions bounds its mass by $CR^{-1/4}$.

Put $z=v'-v$. The side-1 difference has prescribed displacement $z$
in one of its two orientations and costs $D_n(z)$. On side 2 the outer
crossing is free, so the difference costs at most
$C(1+|m\cdot z|)^{-1/4}$. If the side-1 heights tie, its unique local
crossing forces $v=v'$; thus the zero-height convention for $D_n$ is
the correct one. The side-2 outer bridge costs $Cs^{-1/4}$. Finally,
if $v_{\rm out}$ is the earlier side-1 translation, the side-1 outer
bridge costs $D_n(v_{\rm out})$, with
\[
 \sum_{v_{\rm out}\in\mathcal V_s}D_n(v_{\rm out})
 \le C\sum_{h\le Cs}B_h\le Cs^{3/4}.
\]
These upper bounds may be multiplied after dropping avoidance constraints
between the five portions. Lemma~\ref{c:mixed-displacement} sums the two
difference factors. The complete accounting is
\[
 \int N_s^2\,d\mu_a
 \le C\underbrace{s^{3/4}}_{\substack{\text{side-1}\\\text{outer sum}}}
       \qquad\underbrace{s^{1/2}}_{\substack{\text{two}\\\text{differences}}}
       \qquad\underbrace{R^{-1/4}}_{\substack{\text{central}\\\text{chord}}}
       \qquad\underbrace{s^{-1/4}}_{\substack{\text{side-2}\\\text{outer}}}
 =CsR^{-1/4}.
\]
The argument used the two supports and endpoint clearance, not a bound
on the diameter of the input chord.

\paragraph{Both outer landings free.}
Take $s=cR$ and sum over order $R^2$ translations in a fixed positive
sector. Both outer bridges now have one free endpoint, so the first
moment is at least $cR^2R^{-1/2}E$. The same two-list recovery keeps
the inner side-1 crossing prescribed. Its outer bridge is now summed
over its original-support landing, replacing the previous outer sum by
$O(R^2R^{-1/4})$. The other three factors are unchanged. Thus
\[
 \int(N_R^{\rm free})^2\,d\mu_{\rm free}
 \le CR^{2-1/4}R^{1/2}R^{-1/4}R^{-1/4}=CR^{7/4}.
\]
The localized attachments and $O(R)$ translations of bounded input
chords all fit in one convex pure truncation of diameter $O(R)$.
This proves \eqref{c:free-attachment-moments}.
\end{proof}

\subsection{Proof of the interval length law}

\begin{proof}[Proof of Theorem~\ref{c:window-law}]
It suffices to prove the length bounds for $0<\epsilon<1$.
The normalizing mass is comparable to $R^{-1/4}$, as shown above.
If the two inward normals are opposite, the chords are a restriction
of height-$h$ bridges with $h\asymp R$, retaining a fixed fraction of
$B_h$. Lemma~\ref{c:fixed-height-window} gives the result. Otherwise
use Proposition~\ref{c:attachment-proposition}.

For an input event of mass $E=p_0R^{-1/4}$, its fixed-source outputs
have mass at least
\[
 \mu_a(N_s>0)\ge
 \frac{(\int N_s\,d\mu_a)^2}{\int N_s^2\,d\mu_a}
 \ge cp_0^2R^{-1/4}.
\]
All these outputs have endpoints separated by order $R$, so their
total mass is at most $CR^{-1/4}$ by the diameter bound.

Apply this to input lengths below $R^{4/3-\epsilon}$. Choose any fixed
number $M$ of distinct exponents $\gamma_j$ in
$(1-\epsilon/3,1)$ and use $s_j=R^{\gamma_j}$. Choose $\delta>0$
small enough that the $M$ output length windows are disjoint and each
attached length dominates the original short input. The lower mass
constant is independent of this fixed $\delta$. Consequently
$Mcp_0^2\le C+o(1)$. Let $R\to\infty$ first and then $M\to\infty$.
This proves $p_0\to0$ for the short-length event.

For input lengths above $R^{4/3+\epsilon}$, use both free landings.
Equation~\eqref{c:free-attachment-moments} gives output mass at least
\[
 c\frac{R^3E^2}{R^{7/4}}=cp_0^2R^{3/4}.
\]
Every output remains long. Proposition~\ref{c:finite}(iv) bounds its
mass in the encompassing convex domain by
\[
 R^{-4/3-\epsilon}R^{25/12+o(1)}=o(R^{3/4}).
\]
Thus $p_0\to0$ also for the long-length event. All bounds are uniform
over the permitted source interval. Summing exceptional masses and
normalizers over that interval proves the source-summed assertion.
\end{proof}
